\scanpage{308}
\mspage{258}

Aujourd'hui, j'ai \unsure{envie encore|aussi envie|très envie} de donner le \unsure{sorite} des carrés $W$-cartésiens, dans une

\uline{Proposition 19}\,: Soit
\diagram{\begin{tikzcd}[ampersand replacement=\&,row sep=small,column sep=small]
X_1\arrow[d,"{p_1}"']\&X_3\arrow[l,"{f_1}"']\arrow[d,"{f_2}"]\\
S\&X_2\arrow[l,"{p_2}"]
\end{tikzcd}}
un carré comm\unsure{:|.|. $Q$} dans $M$. Considérons des factorisations de $p_1,p_2$ en
\[
X_1\mathrel{\overset{i_1}{\unsure{-|\to}}}\overline X_1\mathrel{\overset{\bar p_1}{\unsure{-|\to}}}S
\qquad X_2\mathrel{\overset{i_2}{\unsure{-|\to}}}\overline X_2\mathrel{\overset{\bar p_2}{\unsure{\leftarrow|-|\to}}}S
\quad\left|\begin{aligned}i_1,i_2&\in W\\ \bar p_1,\bar p_2&\in F_W\end{aligned}\right.
\]
et considérons les flèches induites par $X_3\to X_1\times_S X_2$ (\unsure{déduite} de $Q$) via
\diagram{\begin{tikzcd}[ampersand replacement=\&,row sep=small]
\&\overline X_1\times_S X_2\\
X_1\times_S X_2\arrow[ur]\arrow[r]\arrow[dr]\&\overline X_1\times_S\overline X_2\\
\&X_1\times_S\overline X_2
\end{tikzcd},}
de sorte qu'on trouve \deleted{trois flèches} un diagramme comm\unsure{.}
\diagram{\begin{tikzcd}[ampersand replacement=\&,row sep=small]
\&\overline X_1\times_S X_2\arrow[d,"{j_2}"']\&\deleted{\unsure{\xi}}\\
X_3\arrow[ur,"{\pi_1}"]\arrow[r,"{\pi_0}"]\arrow[dr,"{\pi_2}"']\&\overline X_1\times_S\overline X_2\&\\
\&X_1\times\overline X_2\arrow[u,"{j_1}"]\&
\end{tikzcd},}
où $j_1=i_1\times_S\mathrm{id}_{\overline X_2}$ et $j_2=\mathrm{id}_{\overline X_1}\mathbin{\times}i_2$ sont dans $W$ (car $i_1$ et $i_2$ le sont, et $\overline X_1$ et $\overline X_2$ sont dans $\underline{\mathrm{Fib}}_W S$).

On a donc
\[
(\pi_1\in W)\Longleftrightarrow(\pi_2\in W)\qquad(\unsure{\Longleftrightarrow|\Longrightarrow}\pi_0\in W)
\]
ce qui prouve que la condition $\pi_1\in W$, qui semble dépendre du choix de la factorisation par $i_1$, n'en \uline{dépend pas}, et idem pour $\pi_2\in W$, et \deleted{\unsure{p}}pour $\unsure{\pi|\pi_0}\in W$ (qui semble dépendre à la fois de $i_1$, et de $i_2$).
% archive sheet number: 308
